\subsection{Transmission en ligne indirecte}
\label{sec:transmission_indirecte_sans_retrait}

La figure \ref{fig:sim3} considère le cas d'une transmission en ligne indirecte (par exemple, entre neveux et nièces). Le barème applicable est alors nettement plus sévère : un abattement successoral limité (7 967 euros en 2025) et un taux forfaitaire de 55\% (60\% hors parenté).

\begin{figure}[h!]
    \centering
    \includegraphics[width=\textwidth]{cas_3.png}
    \caption{Différence relative d'héritage net pour un investissement de 100 000 euros avant 70 ans, avec un patrimoine hors enveloppes de 300 000 euros dans le cas d'une transmission entre vers un neveu/nièce.}
    \label{fig:sim3}
\end{figure}

Dans cette configuration, l'avantage de l'assurance-vie est nettement renforcé : l'abattement spécifique de 152 500 euros et les taux réduits (20 à 31,25 \%) contrastent fortement avec le barème applicable au CTO. Même en présence de frais de gestion très élevés (2 \% par an), l'assurance-vie conserve un avantage substantiel dans la plupart des cas, en particulier lorsque les frais d'enveloppe sont optimisés.
